\documentclass[tikz,border=0pt]{standalone}
\usepackage[T1]{fontenc}
\usepackage{helvet}
\usepackage{amsmath,amssymb}
\usepackage{fix-cm}
\renewcommand{\familydefault}{\sfdefault}
\usetikzlibrary{arrows.meta,calc,decorations.pathreplacing,patterns,positioning}
\pdfinfoomitdate=1
\pdftrailerid{}
\definecolor{ink}{HTML}{000000}
\definecolor{blue}{HTML}{3B86FF}
\definecolor{green}{HTML}{00A14B}
\definecolor{red}{HTML}{FC4D30}
\definecolor{purple}{HTML}{A93BFF}
\definecolor{gold}{HTML}{FFC943}
\definecolor{palegold}{HTML}{FFDD92}
\definecolor{muted}{HTML}{747474}
\definecolor{neutral}{HTML}{F1F2F2}
\definecolor{rule}{HTML}{D9D9D9}
\newenvironment{candidate}[1]{%
 \begin{tikzpicture}[x=1cm,y=1cm,line width=.65pt,text=ink,
  font=\fontsize{10}{11.7}\selectfont,
  every node/.style={inner sep=1.4pt},
  flow/.style={-{Stealth[length=1.6mm,width=1.1mm]},line width=.75pt},
  softflow/.style={-{Stealth[length=1.45mm,width=1mm]},line width=.6pt},
  state/.style={circle,draw=blue,fill=blue!8,minimum size=.60cm,line width=.8pt},
  gate/.style={draw=ink,fill=neutral,rounded corners=1.8pt,line width=.65pt},
  observable/.style={draw=ink,fill=gold,rounded corners=1.2pt,line width=.65pt},
  note/.style={text=muted,font=\fontsize{10}{11.7}\selectfont}]
  \useasboundingbox (0,0) rectangle (8.60,5.60);
  \fill[white] (0,0) rectangle (8.60,5.60);
  \node[anchor=west,font=\bfseries\fontsize{12}{13}\selectfont] at (.08,5.37) {#1};
}{\end{tikzpicture}}
\newcommand{\meter}[3]{%
 \begin{scope}[shift={(#1,#2)}]
  \draw[ink,fill=#3,rounded corners=1pt] (-.24,-.22) rectangle (.24,.22);
  \draw[ink,line width=.5pt] (-.16,-.075) arc[start angle=155,end angle=25,radius=.18];
  \draw[ink,line width=.6pt,-{Stealth[length=1.0mm]}] (0,-.075)--(.105,.12);
 \end{scope}}

\begin{document}
\begin{candidate}{c}
% Each group is DEFINED by equality of all declared outputs under every
% allowed continuation, including the empty one. Finite examples do not prove it.
\node at (1.97,4.98) {input-created states};
\node at (6.51,4.98) {future-output behavior};
\draw[blue,fill=blue!5,rounded corners=12pt] (.48,3.04) rectangle (3.55,4.40);
\node[state,minimum size=.61cm] (a1) at (1.11,3.89) {$\rho_a$};
\node[state,draw=red,fill=red!7,minimum size=.61cm] (a2) at (2.12,3.89) {$\sigma_a$};
\fill[ink] (2.99,3.89) circle (.026);
\fill[ink] (3.14,3.89) circle (.026);
\fill[ink] (3.29,3.89) circle (.026);
\node at (1.97,3.26) {one group};
\draw[flow] (3.57,3.73)--(4.68,3.73);
\node[observable,minimum width=3.16cm,minimum height=1.17cm,align=center] at (6.39,3.73) {one behavior A\\for the whole group};
\draw[green,fill=green!5,rounded corners=12pt] (.48,.90) rectangle (3.55,2.26);
\node[state,minimum size=.61cm] (b1) at (1.11,1.75) {$\rho_b$};
\node[state,draw=red,fill=red!7,minimum size=.61cm] (b2) at (2.12,1.75) {$\sigma_b$};
\fill[ink] (2.99,1.75) circle (.026);
\fill[ink] (3.14,1.75) circle (.026);
\fill[ink] (3.29,1.75) circle (.026);
\node at (1.97,1.12) {another group};
\draw[flow] (3.57,1.59)--(4.68,1.59);
\node[observable,minimum width=3.16cm,minimum height=1.17cm,align=center] at (6.39,1.59) {another behavior B\\for the whole group};
\node[text=purple] at (6.38,2.66) {every allowed continuation};
\node[text=green] at (4.28,.35) {within each group, state differences may persist};
\end{candidate}
\end{document}
